
\section{Grundlagen der Mathematik}
\subsection{Aussagenlogik}
\begin{definition}
	Eine Aussage ist eine endliche Zeichenfolge, der genau ein Wahrheitswert, wahr oder falsch, zugeordnet ist.
\end{definition}

Zwei Aussagen $p$ (Hypothese) und $q$ (Konklusion) können wie folgt zu neuen
Aussagen verknüpft werden:
\begin{itemize}
	\item Konjunktion: $p \land q$ (min)
	\item Disjunktion: $p \lor q$ (max)
	\item Negation: $\neg p$
	\item Implikation: $p \implies q$
	\item Äquivalenz: $p \iff q$
\end{itemize}
\paragraph{Prädikatenlogik}
Ein Prädikat ordnet jedem Element $x$ aus einem Universum $X$ eine Aussage
$P(x)$ zu. \note{funktioniert wie Aussage}
\paragraph{Quantoren}
\begin{itemize}
	\item Allquantor $\forall x \in X$
	\item Existenzquantor $\exists x \in X$
\end{itemize}

\subsection{Mengentheoretische Grundlagen}
Mengen bestehen aus Elementen ($x \in M$)
\begin{definition*}
	Naive Definition nach Cantor: Menge Zusammenschluss von bestimmten, wohlunterscheidbaren Objekten zu einem Ganzen.
\end{definition*}
Leere Menge $\emptyset$\\
Beschreibung von Mengen durch Eigenschaften $A:=\{x \in B\mid P(x)\}$

\paragraph{Mengentheoretische Bezeichnungen und Operationen}
Seien $A$ und $B$ Mengen:
\begin{itemize}
	\item Enthaltenseinsbeziehung (Inklusion):\\ $A$ heißt Teilmenge von $B$ ($A
		      \subseteq B$), falls jedes Element von $A$ auch Element von $B$ ist.
	\item echte Teilmenge $A \neq B \land A \subseteq B:A \subset B$
	\item Gleichheit ($A = B$): $A\subseteq B \land B \subseteq A$
	\item Schnitt $A \cap B$\\ $A$ und $B$ disjunkt: $A \cap B = \emptyset$

	\item Vereinigung $A \cup B$
	\item Differenz       $A \setminus B$
	\item Symmetrische Differenz: $A \bigtriangleup B := (A \cup B)\setminus (A\cap B)$
	\item Potenzmenge von $A$: $\mathcal{P}(A) := \{C\mid C\subset A\}$
	\item Kartesisches Produkt: $A \times B := {(a,b)\mid a\in A, b\in B}$ (Tupel)
\end{itemize}

Kardinalität: $|A|$ Anzahl Elemente einer (endlichen) Menge

\subsection{Zermelo-Fraenkel-Axiomatik}
\paragraph{Russelsches Paradoxon}
$M:= \{x\mid x\not\in x\}$ $x=M$\lightning{}
\paragraph{Konstruktiver Aufbau der Mengenlehre}
\begin{enumerate}
	\item Leermengenaxiom: Es gibt eine leere Menge.
	\item Extensionalitätsaxiom: Zwei Mengen sind genau dann gleich, wenn sie dieselben
	      Elemente beinhalten.
	\item Paarmengenaxiom: Für je zwei Mengen $A$ und $B$ existiert eine Menge $\{A,
		      B\}$, die genau $A$ und $B$ als Elemente besitzt.
	\item Vereinigungsaxiom: Für jede Menge $A$ existiert eine Menge, deren Elemente
	      genau die Elemente der Elemente von $A$ sind. (Bsp. Menge enthält Mengen)
	\item Unendlichkeitsaxiom: Es gibt eine Menge $A$, die die leere Menge und für jedes
	      $x \in A$ auch die Menge $x \cup \{x\}$ als Element enthält.
	\item Potenzmengenaxiom: Für jede Menge $A$ existiert eine Menge, deren Elemente
	      genau die Teilmengen von $A$ sind.
	\item Fundierungsaxiom: Jede nichtleere Menge $A$ enthält ein Element $B \in A$
	      derart, dass $A \cap B = \emptyset$. \note{Mengen enthalten sich nicht selbst
		      als Element: $A=\{M\}$}
	\item Informell: Bilder von Mengen unter definierbaren Funktionen sind selbst wieder
	      Mengen.
\end{enumerate}

\subsection{Konstruktion der natürlichen Zahlen}
\begin{definition}
	\begin{align*}
		0:=\emptyset,\quad n^+:=n\cup {n}={0,\dots,n}
	\end{align*}
	In ZF ist Definition $\N:={0,1,2,\dots}$ zulässig.
\end{definition}
Für alle $m,n \in \N$,
\begin{align*}
	m < n \iff m \subset n \iff m \in n
\end{align*}
\paragraph{Addition} rekursiv definiert $\forall m,n \in \N$:
\begin{align*}
	m+0   & := m      \\
	m+n^+ & :=(m+n)^+
\end{align*}
\paragraph{Multiplikation} ebenfalls rekursiv definiert $\forall m,n \in \N$:
\begin{align*}
	m\cdot0    & := 0           \\
	m\cdot n^+ & :=(m\cdot n)+n
\end{align*}
\subsection{Relationen}
\begin{definition}
	Eine Relation zwischen zwei Mengen $A$ und $B$ ist eine Teilmenge $R\subseteq A\times B$.
\end{definition}
Schreibweise: statt $(a,b) \in R$ auch $aRb$ (Infix-Notation)\\
Inverse Relation $R^{-1}\subseteq B \times A$ zu einer Relation $R\subseteq A \times B$ ist definiert als
\begin{align*}
	R^{-1} := \{(b,a)\mid (a,b)\in R\}
\end{align*}
Visualisierung durch Grpahen, etc.\\
Speziallfall $R\subseteq A\times A$ heißt binäre Relation auf $A$.\\
Wichtige Eigenschaften:
\begin{itemize}
	\item reflexiv: $\forall a \in A: (a,a)\in R$
	\item symmetrisch: $\forall a,b \in A: (a,b)\in R \implies (b,a) \in R$
	\item antisymmetrisch: $\forall a,b \in A: ((a,b) \in R \land (b,a)\in R) \implies
		      a=b$
	\item transitiv: $\forall a,b,c \in A: ((a,b)\in R \land (b,c) \in R)\implies
		      (a,c)\in R$
\end{itemize}
\begin{example}
	Bsp.: Teilerrelation, $\leq$-Relation auf $\N$, Diagonalrelation $\Delta_x$, Allrelation $\nabla_x$
\end{example}
\paragraph{Äquivalenzrelation} \
\begin{definition}
	Sei $A$ Menge. Eine Relation $R\subseteq A\times A$ heißt Äquivalenzrelation auf $A$, falls $R$ reflexiv, symmetrisch, transitiv ist.
\end{definition}
Für $x \in A$ heißt $[x]_R:=\{y\in A\mid xRy\}$ die Äquivalenzklasse von $x$ bzgl. $R$. Die Menge $A/R:=\{[x]_R\mid x\in A\}$ aller Äquivalenzklassen bzgl. $R$ wird als Faktormenge von $A$ nach $R$ \note{modulo} bezeichnet.
\begin{theorem}
	Sei $R$ Äquivalenzrelation auf einer Menge $A$. Dann gilt:
	\begin{align*}
		\forall x,y \in A: (x,y) \in R \iff [x]_R=[y]_R
	\end{align*}
\end{theorem}
\prove
\begin{align*}
	\forall x,y \in A: (x,y) \in R \implies [y]_R \subseteq [x]_R \quad \text(*)
\end{align*}
Sei $(x,y)\in R$. Sei $z\in [y]_R$, d.h. $(y,z) \in R$. Wegen Transitivität von $R$: $(x,y)\in R \land (y,z)\in R \implies (x,z) \in R$. Somit $z\in [x]_R$. Somit (*) bewiesen.\\
Sei $(x,y)\in R$. Wegen (*) ist $[y]_R \subseteq [x]_R$. Wegen Symmetrie von $R$ und $(x,y)\in R$ ist auch $(y,x)\in R$, damit nach (*) auch $[x]_R\subseteq [y]_R$. Also $[x]_R=[y]_R$.
Seien $x,y \in A$ mit $[x]_R =[y]_R$. Wegen Reflexivität von $R$ ist $(y,y) \in R$, d.h. $y \in [y]_R$. Daher $y\in [x]_R$ (nach Vorraussetzung), also $(x,y) \in R$. $\Box$
\paragraph{Partitionen}	\
\begin{definition}
	Sei $A$ eine Menge. Eine Partition von $A$ ist eine Menge $\mathcal{C}$ von nicht-leeren Teilmengen von $A$ derart, dass
	\begin{enumerate}
		\item $A=\bigcup \mathcal{C}$, also alle Elemente von $A$ in den Teilmengen
		\item $\forall B,C \in \mathcal{C}: B \neq C \implies B \cap C = \emptyset$, also Teilmengen disjunkt
	\end{enumerate}
\end{definition}
\note{bissl wie Partitionen bei Rechnern}\\
Ist $\mathcal{C}$ eine Partition von $A$, so definiere man eine Relation $Rel_{\mathcal{C}}\subseteq A \times A$ durch
\begin{align*}
	Rel_{\mathcal{C}} := \{(x,y)\in A^2\mid \exists C \in \mathcal{C}: (x\in C)\land(y \in C)\}
\end{align*}
\begin{theorem}
	Sei $A$ Menge.
	\begin{enumerate}
		\item Ist R eine Äquivalenzrelation auf $A$, dann ist $A/R$ eine Partition auf $A$,
		      und es gilt $Rel_{A/R}=R$.
		\item Ist $\mathcal{C}$ eine Partition auf $A$, dann ist $Rel_{\mathcal{C}}$ eine
		      Äquivalenzrelation auf $A$, und es gilt $A/Rel_{\mathcal{C}}=\mathcal{C}$.
	\end{enumerate}
\end{theorem}
\subsection{Abbildung (Funktion)}
\begin{definition*}
	Eine Abbildung $f$ von einer Menge $A$ nach einer Menge $B$ ordnet jedem $x \in A$ genau ein Element aus $B$ zu, welches dann mit $f(x)$ bezeichnet wird.
\end{definition*}
(Schreibweise $f:A \to B, x \mapsto f(x)$)\\
Genauer:
\begin{definition}
	Eine Abbildung $f$ ist ein Tripel $(A,f^{\bullet}, B)$ bestehend aus
	\begin{itemize}
		\item einer Menge $A$ (Definitionsbereich)
		\item einer Menge $B$ (Zielbereich)
		\item einer Relation $f^{\bullet}\subseteq A \times B$ (Graph von $f$
		      ($\mathrm{graph}(f)$))
	\end{itemize}
	mit folgenden Eigenschaften:
	\begin{enumerate}
		\item $\forall a\in A, \exists b\in B: (a,b)\in f^{\bullet}$.
		\item $\forall a \in A, \forall b, b' \in B: ((a,b)\in f^{\bullet} \land (a,b')\in f^{\bullet}) \implies b=b'$
	\end{enumerate}
\end{definition}
Für das zu einem $a \in A$ eindeutig bestimmte $b\in B$ mit $(a,b)\in f^{\bullet}$ schreibt man $f(a):=b$ und nennt $f(a)$ das Bild von $a$ unter $f$. Damit:
\begin{align*}
	f^{\bullet}=\{(a,b)\in A \times B\mid f(a)=b\}
\end{align*}
\paragraph{Spezielle Eigenschaften von Abbildungen}
\begin{itemize}
	\item $f:A\to B$ surjektiv: $\forall b \in B \exists a \in A: f(a)=b$
	\item $f:A\to B$ injektiv: $\forall a, a' \in A: f(a)=f(a')\implies a = a'$
	\item $f:A\to B$ bijektiv: $f$ surjektiv und injektiv
\end{itemize}
\paragraph{Weitere Bezeichnungen} Sei $f:A\to B$ Abbildung. Für $A'\subseteq A$ heißt $f(A'):=\{f(a)\mid a\in
	A'\}=\{b \in B\mid \exists a \in A': b=f(a)\}$ das Bild von $A'$ unter $f$.\\
$im(f):=f(A)$ Bild von $f$.\\ Für $B' \subseteq B$ heißt $f^{-1}(B'):=\{a\in
	A\mid f(a)\in B'\}$ das Urbild von $B'$ unter $f$.\\ Die Einschränkung von $f$
auf auf eine Teilmenge $A'\subseteq A$ ist die Abbildung
\begin{align*}
	f|_{A'}:A'\to B, a \mapsto f(a)
\end{align*}
d.h. $(f|_{A'})^{\bullet}=f^{\bullet} \cap (A'\times B)$. Die binäre Relation
\begin{align*}
	\ker(f):=\{(a,a')\in A \times A\mid f(a)=f(a')\}
\end{align*}
heißt der Kern von $f$, und $\ker(f)$ ist eine Äquivalenzrelation auf $A$.
\begin{example}\
	\begin{itemize}
		\item $A$ Menge. Dann $\mathrm{id}_A: A \to A$ bijektiv und $\mathrm{graph}(\mathrm{id}_A)=\Delta_A$
		\item Seien $A_1, A_2$ Mengen. Die Abb. \begin{align*}
			      p_1 & : A_1 \times A_2 \to A_1 \\
			      p_2 & : A_1 \times A_2 \to A_2
		      \end{align*}
		      heißen Projektionen.
		\item $R$ Äquivalenzrelation auf Menge $A$. Dann ist die Abb.
		      \begin{align*}
			      \pi_R: A \to A/R
		      \end{align*}
		      surjektiv, und es gilt \begin{align*}
			      \ker(\pi_R) & =\{(x,y)\in A^2\mid \pi_R(x)=pi_R(y)\} \\
			                  & =\{(x,y)\in A^2\mid [x]_R=[y]_R\}      \\
			                  & =\{(x,y)\in A^2\mid (x,y)\in R\}=R
		      \end{align*}
	\end{itemize}
\end{example}
\paragraph{Komposition von Abbildungen}
$A,B,C$ Mengen. Für Abb. $f: A \to B$ und $g:B\to C$ heißt
\begin{align*}
	g\circ f: A \to C, x \mapsto g(f(x))
\end{align*}
die Komposition (Verkettung) `$g$ nach $f$'.\\
Für $f: A \to A$ und $n\in \N$ ist
\begin{align*}
	f^n:= f \circ \dots \circ f:A \to A
\end{align*}
Genauer
\begin{align*}
	f^0:=\mathrm{id}_A:A\to A, \\
	f^{n+1}:=f^n\circ f: A \to A\quad(n\in \N)
\end{align*}
\paragraph{Umkehrabbildungen}
Betrachte Abb. $f: A\to B$. Genau dann ist die inverse Relation
\begin{align*}
	(f^{\bullet})^{-1}\subseteq B \times A
\end{align*}
der Graph einer Abb. von $B$ nach $A$, wenn $f$ bijektiv ist. In diesem Fall nennt man $f^{-1}:B\to A$ mit $(f^{-1})^{\bullet}:=(f^{\bullet})^{-1}$, d.h. $f^{-1}(b)=a: f(a)=b$, die Umkehrabbildung (Umkehrfunktion) von $f$.
\subsection{Gleichmächtigkeit}
$A,B$ Mengen. $A$ und $B$ sind gleichmächtig ($|A|=|B|$):$\iff$ $\exists$ bijektive Abb. $f: A \to B$\\
Weiter:
\begin{align*}
	|A|\neq |B| & :\iff \not(|A|=|B|)                     \\
	|A|\leq |B| & :\iff \exists f: A\to B \text{injektiv} \\
	|A|< |B|    & :\iff |A|\neq |B| \land |A|\leq |B|
\end{align*}
\begin{example}
	$\N, \Q, \Z, \N\times \N$ sind gleichmächtig
\end{example}
\begin{theorem}
	(Cantor) Für jede Menge $A$ ist $|A| < |\mathcal{P}(A)|$.
\end{theorem}
\prove 
\begin{itemize}
	\item injektiv: Sei $A$ Menge. Dann ist $A\to \mathcal{P}, a \mapsto \{a\}$ injektiv. Also $|A|\leq |\mathcal{P}(A)|$
	\item nicht bijektiv: Angenommen, $|A|=|\mathcal{P}(A)|$. Insbesondere gibt es eine surjektive Abbildung $f: A \to \mathcal{P}(A)$. Betrachte $B:=\{x\in A\mid x \not\in f(x)\} \in \mathcal{P}(A)$. Da $f$ surjektiv, existiert $b \in A$ mit $f(b)=B$. Nun entweder $b\in B$ oder $b\not\in B$ \lightning. Somit $|A|\neq|\mathcal{P}(A)|$ und daher $|A| < |\mathcal{P}(A)|$.
\end{itemize}
\begin{theorem}[Bernstein-Schröder]
	Für beliebige Mengen $A$ und $B$ gilt: Ist $|A|\leq |B|$ und $|B|\leq |A|$, dann $|A|=|B|$.
\end{theorem}
\prove Siehe Bücher
\paragraph{Bemerkung} Sei $f: A\to B$ eine Abbildung.
\begin{enumerate}
	\item $f$ bijektiv $\iff \exists g: B\to A: g\circ f=\mathrm{id}_A, f\circ g=\mathrm{id}_B$. In diesem Fall ist $g=f^{-1}$.
	\item Falls $A \neq \emptyset$, dann: $f$ injektiv $\iff \exists g:B\to A: g\circ f=\mathrm{id}_A$
\end{enumerate}
\subsection{Auswahlaxiom}
Eine zum Vorherigen analoge Beschreibung surjektiver Abbildungen (die man in der ZF-Axiomatik nicht beweisen kann) bildet den Kern eines weiteren mengentheoretischen Axioms:\\
Für jede surjektive Abbildung $f: A\to B$ existiert eine Abbildung $g: B\to A$ mit $f\circ g=\mathrm{id}_B$ 
\subsection{Beweisprinzipien}
$A\implies B$
\begin{itemize}
	\item Direkter Beweis $A\implies B$
	\item Indirekter Beweis $((\neg B)\implies (\neg A))$
\end{itemize}
Ringschluss, Vollständige Induktion
\section{Körper}
\subsection{Gruppe}
\begin{definition}[Gruppe]
	Ein Paar $(G, \cdot)$ bestehend aus einer Menge $G$ und einer Abbildung $\cdot: G\times G\to G$ heißt Gruppe, falls folgende Axiome erfüllt sind:
	\begin{enumerateroman}
	\item Assoziativität
	\item Neutrales Element
	\item Inverses Element
	\end{enumerateroman}
\end{definition}
\begin{definition}[abelsche Gruppe]
	Gruppe $(G, \cdot)$ heißt abelsch, falls zusätzlich Kommutativität
\end{definition}
Sei $X$ Menge und $\mathrm{Sym}(X):=\{f\in X^x \mid f \text{ bijektiv}\}$, d.h. die Menge aller bijektiven Abbildungen von $X$ nach $X$. Dann ist $(\mathrm{Sym}(X), \cdot)$ eine Gruppe, die sogenannte (volle) symmetrische Gruppe über $X$. 
\subsection{Körper}
\begin{definition}
	Ein Tripel $\K$ $(K, +, \cdot)$ mit
	\begin{itemize}
		\item $+$: $K\times K\to K, (x,y)\mapsto x+y$
		\item $\cdot$: $K\times K\to K, (x,y)\mapsto xy$
	\end{itemize}
	derart, dass Folgende Eigenschaften gelten:
	\begin{itemize}
		\item $(K,+)$ abelsche Gruppe mit Nullelement $0_{K}$
		\item $(K\setminus 0, \cdot)$ abelsche Gruppe mit Einselelement $1_{K}$
		\item $0_{K}\neq 1_{K}$ 
		\item Distributivität
	\end{itemize}
\end{definition}
\paragraph{Folgerungen}
\begin{itemize}
	\item $0\cdot a=0$
	\item $(-1)a=-a$
	\item $(-a)(-b)=ab$
	\item $ab=0\iff a=0\lor b=0$
\end{itemize}
\subsection{Restklassenarithmetik}
\begin{definition}
	Sei $n\in \N^+$:\\
	$\equiv_n$ definiert auf $\Z$ eine Äquivalenzrelation 
	\begin{align*}
		x \equiv_n y \iff x \equiv y\ \mod n \iff n|(x-y)
	\end{align*}
	Für $x\in \Z$ heißt die Äquivalenzklasse
	\begin{align*}
		[x]_{\mod n}:=\{y\in \Z\mid x\equiv y\ \mod n\}
	\end{align*}
	die Restklasse (oder auch Kongruenzklasse) von von $x$ modulo $n$.\\[0.2cm]
	Menge aller Restklassen: $\Z/(\mod n)$
\end{definition}
\begin{itemize}
	\item Addition: $[x]\oplus\ [y]:=[x+y]$
	\item Multiplikation: $[x]\odot\ [y]:=[x\cdot y]$
\end{itemize}
\paragraph{Bemerkung} Man musss beweisen, dass diese Operationen `wohldefiniert' sind $\iff$ Zuordnungen sind nicht abhängig von der Auswahl der Reprösentanten der jeweiligen Kongruenzklasse!
\paragraph{Alternative Methode} Teilen mit Rest \begin{align*}
	\forall x \in \Z \exists!(k,r)\in \Z\times \Z_n: x = kn+r 
\end{align*}
Bezeichnung: $x\ \mod n := r$ 
\begin{align*}
	\forall x,y \in \Z: x\equiv y \mod n \iff x \mod n = y \mod n
\end{align*}
\paragraph{Addition und Multiplikation auf $\Z_n$}
\begin{align*}
	x\oplus y&:= (x+y\mod n)\\
	x\odot y&:= (x\cdot y\mod n)
\end{align*}
Die Abbildung $\phi: \Z_n\to \Z/(\mod n), x \mapsto [x]_{\mod n}$ ist eine Bijektion für $x \in \Z$.\\
$\forall x,y \in \Z_n$:
\begin{align*}
	\phi(x+y)&=\phi(x)\oplus\phi(y)\\
	\phi(x\cdot y)&=\phi(x)\odot\phi(y)
\end{align*}
$\phi$ $(\Z_n, +, \cdot)$ $\to$ $\phi$ $(\Z_{\mod n}, \oplus, \odot)$ ist ein Isomorphismus
\begin{theorem}
	Sei $n\in\N^+$. Dann:
	\begin{align*}
		(\Z_n,+,\cdot) \text{ist Körper} \iff (\Z_{\mod n}, \oplus, \odot) \text{ist Körper} \iff n\in\P
	\end{align*}
\end{theorem}
\paragraph{Bemerkung}
\begin{itemize}
	\item $\forall n \in \N^+$ sind $(\Z_n,+,\cdot)$ Ringe (nur Halbgruppe bzgl. $\cdot$ (Assoziativität und abgeschlossen))
	\item $\forall p \in \P, \forall m \in \N^+ \exists! \K: |\K|=p^m$
\end{itemize}
\subsection{Der Körper der Komplexen Zahlen}
Konstruktion mittels Tupel $(x,y)\in \R^2$ sowie den Verknüpfungen $\odot$ und $\oplus$ sei dem Leser überlassen.
\begin{align*}
	\C := \{a+bi \mid a,b\in \R\}
\end{align*}
Die Menge $\C$ bildet mit den vorangegangenen Verknüpfung $\oplus$ und $\odot$ den `Körper der komplexen Zahlen' $(\C, \oplus, \odot)$.\\
Operationen in $\C$:
\begin{itemize}
	\item $\Re(z)=a$ Realteil
	\item $\Im(z)=b$ Imaginärteil
\end{itemize}
Jede reelle Zahl $a\in\R$ ist als komplexe $a+0i\in\C$ zu versthen. Darum schreibt man $\R \subset \C$. Das heißt man identifiziert die Menge $\R$ mit ihrem Bild unter der Injektion $\R\to\C,a\mapsto a+bi$.
\section{Vektorräume} 
\begin{definition}
	Jeden Vektroräume liegt ein Körper $\K$ $(K,+,\cdot)$ zugrunde. Die Elemente von $\K$ heißen oft Skalare.
\end{definition}
\subsection{Vektorraum}
\begin{definition}
	Sei $\K$ Körper. Vektorraum über $\K$ (Kurz $\K$-Vektorraum) ist eine Menge $V$ mit zwei Abbildungen
	\begin{itemize}
		\item Addition: $V\times V \to V, (u,v)\mapsto u+v$
		\item Skalarmultiplikation: $\K\times V \to V, (\lambda, v)\mapsto \lambda v$
	\end{itemize}
	derart, dass folgende Axiome erfüllt sind:
	\begin{enumerateroman}
		\item $(V,+)$ abelsche Gruppe
		\item $\forall \lambda, \mu \in \K \forall v\in V: (\lambda \cdot \mu)\cdot v = \lambda \cdot (\mu \cdot v)$
		\item $\forall v \in V:1v=v$
		\item $\forall \lambda \in \K \forall u,v \in V: \lambda \cdot (u+v) = \lambda u + \lambda v$
		\item $\forall \lambda, \mu \in \K \forall v \in V: (\lambda+\mu)\cdot v = \lambda v + \mu v$
	\end{enumerateroman}
	Vektor: Element von $V$
\end{definition}
\begin{example}[Vektorraum $\K^n$]
Sei $\K$ Körper und $n\in \N$. Die Menge $\K^n$ aller $n$-Tupel bildet einen $\K$-VR.	
\end{example}
\begin{example}[Vektorräume aus Körpererweiterungen]
	Sei $\K$ Körper. $L\subseteq \K$ heißt Unterkörper (auch Teilkörper) von $\K$. Man schreibt dann $L \leq \K$, sofern gilt:
	\begin{itemize}
		\item $0_{\K}, 1_{\K} \in L$
		\item $\forall a,b \in L: a+b \in L$
		\item $\forall a \in L: -a\in L$
		\item $\forall a, b \in L: a\cdot b \in L$
		\item $\forall a\in L\setminus\{0\}: a^{-1}\in L$
	\end{itemize}
	\note{abgeschlossen, inverses und neutrales Element}\\
	Sei $L\leq \K$, dann $\K$-Vektorraum über $L$.
\end{example}
\subsection{Funktionenräume}
\begin{definition}
	Sei $\K$ Körper und $X$ Menge. Dann bildet $\K^X := \{f\mid f:X\to \K\}$ einen $\K$-VR mit folgenden Operationen:
	\begin{itemize}
		\item Addition: $(f+g)(x)=f(x)+g(x)$
		\item Skalarmultiplikation $(\lambda f)(x)=\lambda f(x)$
		\item Nullvektor $x\mapsto 0_{\K}$
	\end{itemize}
\end{definition}
\begin{example}[Potenzmenge als $\Z$-Vektorraum]
	Sei $A$ Menge. Dann bildet $\mathcal{P}(A)$ einen $\Z_2$-VR mit folgenden Operationen:
	\begin{itemize}
		\item Addition: $X+Y=(X\cup Y)\setminus(X\cap Y)$
		\item Skalarmultiplikation
		\item Nullvektor
		\item Inverses bzgl. Addition
	\end{itemize}
\end{example}
\begin{lemma}
	Sei $V$ $\K$-VR, $v \in V$. Dann gilt:
	\begin{enumerateroman}
		\item $0v=0$
		\item $(-1)v=-v$
		\item $\lambda \in \K: \lambda v=0 \iff (\lambda=0)\lor (v=0)$
	\end{enumerateroman}
\end{lemma}
\prove
\begin{enumerateroman}
	\item \begin{align*}
		0&= 0v-(0v)\\
		&=(0+0)v-(0v)\\
		&= 0v+0v-0v\\
		&=0v
	\end{align*}
	\item \begin{align*}
		0&=(1-1)v=v+(-1)v\\
		-v=(-1)v
	\end{align*}
\end{enumerateroman}
\begin{definition}
	$V$ $\K$-VR. Eine Teilmenge $U\subseteq V$ heißt Untervektorraum (UVR) von $V$ und man schreibt $U\leq V$, wenn gilt:
	\begin{enumerate}
		\item $U\neq \emptyset$
		\item $\forall u,v \in U: u+v \in U$
		\item $\forall \lambda \in \K \forall u \in U: \lambda u \in U$
	\end{enumerate}
\end{definition}
\begin{lemma}
	$V$ $\K$-VR, $U\leq V$. Dann gilt:
	\begin{enumerate}
		\item $0_v\in U$
		\item $\forall u \in U: -u \in U$
		\item Skalarmutliplikation und Addition in $V$
	\end{enumerate}
\end{lemma}
\prove Sei dem Leser überlassen
\paragraph{Bemerkung} Für jeden Vektorraum $V$ (über einem beliebigen Körper) ist $V\leq V$ und $\{0\}\leq V$.
\begin{example}
	Sei $\K$ Körper.
	\begin{enumerate}
		\item Sei $n\in \N$ und $M\subseteq \{1,\dots, n\}$. Dann
		\begin{align*}
			\left\{\begin{array}{c}
				a_1\\\dots\\a_n
			\end{array}\mid \forall i \in M: a_i = 0\right\} \leq \K^n
		\end{align*}
		\item Sei $X$ beliebige Menge. Dann bildet die Menge $\K^X$ aller endlich getragenen Funktionen von $X$ nach $\K$ einen UVR des $\K$-VR $\K^X$
	\end{enumerate}
\end{example}
\begin{lemma}
	$V$ $\K$-VR. Sind $U_1$ und $U_2$ UVR von $V$, dann ist auch $U_1\cap U_2\leq V$.
\end{lemma}
\prove Sei dem Leser überlassen.
\paragraph{Bemerkung} Allgemein gilt: Ist $M$ eine Menge von Untervektorräumen von $V$, dann ist auch $\bigcap M$ UVR von $V$.\\
Unter Vereinigung kein UVR gegeben.
\begin{definition}
	$V$ $\K$-VR.
\begin{enumerate}
	\item 
	\item 
\end{enumerate}
	\begin{enumerate}
		\item Für $S, T\subseteq V$ sei $S+T:=\{u+v\mid u\in S, v\in T\}$ FÜr $v\in V$ sei $Kv:=\{\lambda v\mid \lambda \in \K\}$. Für $m, n\in \N, m\leq n, v_m , \dots , v_n \in V$ definiert man die Summe $\sum\limits_{i=m}^n v_i \in V$ mittels Rekursion. Für $m, n \in \N$ mit $m>n$ sei $\sum\limits_{i=m}^n v_i=0$
		\item Für $n\in\N$, $v_1,\dots,v_n\in V$ und $\lambda_1,\dots,\lambda_n\in\K$ heißt
		\begin{align*}
			\sum\limits_{i=1}^n \lambda_iv_i= \lambda_1v_1+\dots+\lambda_nv_n
		\end{align*}
		eine Linearkombination der Vektoren. Für $S\subseteq V$ heißt $\mathrm{Lin}_V(S):=\mathrm{Lin}_V(S):=\{\sum\limits_{i=1}^n\lambda_iv_i\mid n \in \N, \lambda \in \K, v\in S\}$ die lineare Hülle von $S$ in $V$.\\ Spezialfall: $\mathrm{Lin}_V(\emptyset)=\{0\}$.
		\item Eine Teilmenge $S \subseteq V$ heißt Erzeugendensystem des $\K$-VR $V$, wenn $V=\mathrm{Lin}_V(S)$.
	\end{enumerate}
\end{definition}
\begin{theorem}
	Sei $V$ $\K$-VR und $S\subseteq V$. Dann ist $\mathrm{Lin}_V(S)$ der kleinste $UVR$ von $V$, der S als Teilmenge enthält, d.h. 
	\begin{enumerate}
		\item $S\subseteq \mathrm{Lin}_V(S)\leq V$
		\item $\forall U \leq V: S \subseteq U \implies \mathrm{Lin}_V(S) \subseteq U$
	\end{enumerate}
\end{theorem}
\prove Sei dem Leser überlassen.
\paragraph{Bemerkung} $V$ $\K$-VR. Dann gilt:
\begin{itemize}
	\item Die Abbildung $\mathcal{P}\to\mathcal{P}, S \mapsto \mathrm{Lin}_V(S)$ ist ein Hüllenoperator, d.h. $\forall S,T\in \mathcal{P}$ gilt:
	\begin{enumerate}
		\item $S\subseteq \mathrm{Lin}_V(S)$ (Extensivität)
		\item $S\subseteq T \implies \mathrm{Lin}_V(S)\subseteq\mathrm{Lin}_V(T)$ (Monotonie)
		\item $\mathrm{Lin}_V(\mathrm{Lin}_V(S))=\mathrm{Lin}_V(S)$ (Idempotenz)
	\end{enumerate}
	\item Für jedes $S\subseteq V$ ist
	\begin{align*}
		\mathrm{Lin}_V(S)=\bigcap \{U\in \mathcal{P}(V)\mid S\subseteq U \leq V\}
	\end{align*}
\end{itemize}
\prove 
\begin{itemize}
	\item (i) und (ii) nach oben. (iii): Sei $S\subseteq V$. Nach (i) ist $S\subseteq \mathrm{Lin}_V(S)$, also auch $\mathrm{Lin}_V(S)\subseteq \mathrm{Lin}_V(\mathrm{Lin}_V(S))$. Weiter ist $\mathrm{Lin}_V(S)\leq V$. Daher $\mathrm{Lin}_V(\mathrm{Lin}_V(S))\subseteq \mathrm{Lin}_V(S)$.
	\item Sei $S \subseteq V$ und $W:=\bigcap \{U\in \mathcal{P}(V)\mid S\subseteq U\leq V\}$. Zu zeigen $\mathrm{Lin}_V(S) = W$. $W\leq V$. Wegen $S\subseteq W$ also $\mathrm{Lin}_V(S)\subseteq W$. Da $S\subseteq \mathrm{Lin}_V(S)\leq V$, gilt auch $W\subseteq \mathrm{Lin}_V(S)$ 
\end{itemize}
\paragraph{Proposition} Sei $V$ $\K$-VR. Für alle $n\in N$ und $v_i\in V$ gilt:
\begin{align*}
	\mathrm{Lin}_V(v_i)=\sum\K v_i
\end{align*}
\begin{definition}
	$V$ $\K$-VR. Ein $n$-Tupel $(v_1,\dots, v_n)\in V^n$ heißt linear unabhängig, falls
	\begin{align*}
		\forall \lambda_1, \dots, \lambda_n \in \K: \lambda_1v_1+\dots +\lambda_nv_n=0\implies \lambda_1=\dots=\lambda_n =0
	\end{align*}
	und $(v_1, \dots, v_n)$ heißt linear abhängig, falls nicht linear unabhängig. Eine Teilmenge $U\subseteq V$ heißt linear unabhänig, wenn jedes endliche Tupel paarweise verschiedener Elemeente aus $U$ linear unabhängig ist, d.h. $\forall n \in N \forall (v_i, \dots, v_n)\in U^n$:
	\begin{align*}
		\left(\forall i,j\in \{1,\dots,n\}: i\neq j\implies v_i\neq v_j\right) \implies \left((v_1,\dots, v_n) \text{ linear unabhängig}\right)
	\end{align*}
\end{definition}
\paragraph{Bemerkung}
\begin{itemize}
	\item Ein Vektor $v\in V$ ist genau dann linear abhängig, wenn $v=0$.
	\item Ein Tupel $(v_1,\dots, v_n)\in V^n$ mit $n\geq 2$ ist genau dann linear abhängig, wenn
	\begin{align*}
		\exists i \in \{1,\dots,n\} \exists \lambda_1,\dots,\lambda_{i-1},\lambda_{i+1},\dots, \lambda_n\in \K:\\ v_i = \lambda_1v_1+\dots + \lambda_{i-1}v_{i-1}+\lambda_{i+1}v_{i+1}+\dots + \dots \lambda_nv_n
	\end{align*}
	\item Ein Tupel $(v_1, \dots, v_n)\in V^n$ ist genau dann linear unabhängig, wenn $v_i$ paarweise verschieden und $\{v_1,\dots,v_n\}$ linear unabhängig.
	\item Seien $U\subseteq \tilde{U}\subseteq V$. Dann gilt $\tilde{U}$ linear abhängig $\implies$ $U$ linear unabhängig. 
\end{itemize}
\begin{example}
	Im $\R$-VR $\R^2$:
	\begin{itemize}
		\item $((0,1)^T, (1,0)^T)$ linear unabhängig.
		\item $((1,2)^T, (2,4)^T)$ linear abhängig.
		\item $((1,2)^T, (\pi, 2\pi)^T)$ linear abhängig im $\R$-VR $\R^2$, aber linear unabhängig im $Q$-VR $\R^2$.
		\item Je drei Vektoren im $\R$-VR $\R^2$ sind linear abhängig
	\end{itemize}
\end{example}
\section{Basen und Dimensionen}
\begin{definition}
	$V$ $K$-VR. Eine Teilmenge $B\subseteq V$ heißt Basis von $V$:
	\begin{enumerateroman}
	\item $B$ is linear unabhängig
	\item $B$ ist Erzeugendensystem von $V$ ($\mathrm{Lin}_V(B)=V$)
	\end{enumerateroman}
	Ein Tupel $B=(v_1, \dots, v_n)\in V^n$ mit $n\in \N$ heißt geordnete Basis von $V$:
	\begin{enumerate}
		\item $v_1, \dots, v_n$ paarweise verschieden
		\item $\{v_1, \dots, v_n\}$ Basis von $V$
	\end{enumerate}
\end{definition}
\begin{theorem}[Eindeutigkeit der Basiskoordinaten]
	$V$ $K$-VR, $B=(v_1,\dots, v_n)$ geordnete Basis von $V$. Dann:
	\begin{enumerate}
		\item Für jedes $v \in V$ existiert genau ein $n$-Tupel $(\lambda_1, \dots, \lambda_n)\in K^n$ mit $v=\sum\limits_{i=1}^n \lambda_iv_i$
		\item Die Abbildung \begin{align*}
			\Phi: K^n \to V, (\lambda_1,\dots, \lambda_n)\mapsto \lambda_1v_1+\dots + \lambda_nv_n
		\end{align*}
		ist eine Bijektion und heißt der kanonische Basisisomorphismus bezüglich $B$.
	\end{enumerate}
\end{theorem}
\prove Sei $v\in V$. Da $\{v_1,\dots, v_n\}$ Erzeugendensystem von $V$, ist $V=\mathrm{Lin}_V(v_1,\dots,v_n)=Kv_1+\dots+Kv_n$. Also existieren $\lambda_1,\dots, \lambda_n \in K$ mit $v=\lambda_1v_1+\dots +\lambda_nv_n$.\\
Rest sei dem Leser überlassen.
\paragraph{Bemerkung} $K$ Körper, $n\in \N$. Dann $B=(e_1,\dots, e_n)$ mit Einheitsvektoren eine Basis des $K$-VRs $K^n$, die sogenannte Standardbasis. Für jedes $a=(\lambda_1,\dots,\lambda_n) \in K^n$ ist $a\in \mathrm{Lin}_{K^n}(e_1,\dots, e_n)$
\begin{theorem}[Charakterisierungssatz für Basen]
	$V$ $K$-VR und $B\subseteq V$. Dann gilt:
	\begin{enumerate}
		\item ``Basisergänzungssatz'': Genau dann ist $B$ eine Basis von $V$, wenn $B$ maximale linear unabhängige Teilmenge von $V$ ist, d.h. $B$ linear unabhänig und $\forall \tilde{B}\subseteq V: B \not\subseteq\tilde{B}\implies \tilde{B}$ linear abhängig.
		\item ``Basisauswahlsatz'': Genau dann ist $B$ Basis von $V$, wenn $B$ minimales Erzeugendensystem von von $V$.
	\end{enumerate}
\end{theorem}
\prove \ \\Sei $B$ Basis von $V$. Dann $B$ linear unabhängig. Sei $\tilde{B}\subseteq V$ mit $B\not\subseteq \tilde{B}$. Dann gibt es $v\in\tilde{B}\setminus B$. Da $\mathrm{Lin}_V(B)=V$, existieren $n\in \N$, $\lambda_i\in K$ und paarweise verschiedene $v_i\in B$ mit $\sum \lambda_iv_i= v$. Damit $v_1, \dots, v_n, v \in \tilde{B}$ paarweise verschieden, $\sum \lambda_iv_i+(-1)v=0$, also $\tilde{B}$ linear abhängig.\\
Sei $B$ maximal linear unabhängig in $V$. Zu zeigen: $B$ Basis von $V$. Widerspruchsbeweis: Angenommen, $B$ keine Basis von $V$, d.h. $\mathrm{Lin}_V(B)\neq V$. Dann existiert $v\in V\setminus\mathrm{Lin}_V(B)$. Zeige $\tilde{B}:=B\cup \{v\}$ linear unabhängig. Seien $n\in \N$,$\lambda_1, \dots, \lambda_{i+1}\in K$ und $v_1, \dots, v_n \in B$ derart, dass $v_i$ paarweise verschieden und $\lambda_1v_1+\dots+\lambda_nv_n+\lambda_{n+1}v=0$. Wäre $\lambda_{n+1}\neq 0$, dann wäre $v$ Linearkombination der $v_i$ \lightning. Also $\lambda_{n+1}=0$ und daher $\tilde{B}$ linear unabhängig. Wenn $\tilde{B}$ linear unabhängig ist widerspricht das maximaler linearer Unabhängigkeit von $B$ \lightning
\begin{theorem}
	Jeder Vektorraum hat eine Basis.
\end{theorem}
\prove Sei $V$ $K$-VR.\\
Speziallfall: Ist $V$ endlich erzeugt, d.h. $\exists S \subseteq V:$ $S$ endlich $\land$ $V=\mathrm{Lin}_V(S)$, dann enthält $S$ ein minimales Erzeugendenystem $B$ von $V$. Dann ist $B$ Basis von $V$. Allgemeiner Fall (nicht notwendigerweise endlich) benötigt Lemma von Zorn (äquivalent zum Auswahlaxiom).
\paragraph{Bemerkung} $K$ endlicher Körper, $q:=|K|$. Hat ein $K$-VR $V$ eine Basis mit $n$ ELementen, dann  $|V|=|K^n|=|K|^n=q^n$, also hat jede Basis von $V$ genau $n$ Elemente. DIese Beobachtung lässt sich auf Vektorräume über beliebigen Körpern verallgmeinern.
\begin{lemma}[Austauschlemma]
	$V$ $K$-VR. Sei $(v_1,\dots,v_n)$ Basis von $V$, $\lambda_1, \dots, \lambda_n \in K$ und $w=\lambda_1v_1+\dots+\lambda_nv_n$. Dann lässt sich ein Element der Basis für das $\lambda_j\neq 0$ $v_j$ mit $w$ austauschen.
\end{lemma}
\prove Sei dem Leser überlassen
\begin{theorem}[Austauschsatz von Steinitz]
	Sei $V$ $K$-VR. Sei $B\subseteq V$ endliche Basis von $V$ und $n:=|B|$. Sei $C\subseteq V$ endlich und linear unabhängig, $m:=|C|$. Dann:
	\begin{enumerateroman}
	\item $|C|\leq |B|$, d. h. $m\leq n$
	\item Es existieren paarweise verschiedene Vektoren $b_1, \dots, b_{n-m}\in B\setminus C$ derart, dass $C \cup {b_1, \dots, b_{n-m}}$ Basis von $V$ ist.
	\end{enumerateroman}
\end{theorem}
\prove induktiv 
\paragraph{Folgerung} Sei $V$ $K$-VR, $B$ endliche Basis von $V$, $n:=|B|$.\begin{enumerate}
	\item Alle Basen von $V$ haben gleiche Kardinalität
	\item Jede linear unabhängige Teilmenge $C\subseteq V$ mit $|C|=n$ ist Basis von $V$.
	\item Ist $U\leq V$, so hat jede Basis von $U$ höchstens $n$ Elemente und kann zu einer Basis von $V$ ergänzt werden.
\end{enumerate}
\begin{definition}
		Sei $VV$ eine $K$-VR. Besitzt $V$ endliche Basis $B\subseteq V$, dann nennt man $\dim_K(V):=|B|$ die Dimension des $K$-VRs $V$. Falls $V$ keine endliche Basis besitzt, dann $\dim_K(V):=\infty$.
\end{definition}
\paragraph{Bemerkung} $V$ $K$-VR. Dann
\begin{enumerate}
	\item $\dim_K(V)$ ist wohldefiniert, d.h. Definition hängt nicht von der Wahl der Basis ab.
	\item $\dim_K(V)=0\iff V=\{0\}$
	\item Ist $U\leq V$, dann $\dim_K(U)\leq \dim_K(V)$ und für endliche Dimension $\dim_K(U)=\dim_K(V)\implies U=V$
\end{enumerate}
\begin{example}
	$K$ Körper. Für jedes $n\in\N$ ist
	\begin{align*}
		\dim(K^n)=n
	\end{align*}
	denn bspw. die Standardbasis hat Mächtigkeit $n$
\end{example}
\begin{definition}
	$V$ $K$-VR, $U_1, U_2 \leq V$. Dann heißt 
	\begin{align*}
		U_1+U_2=\{u+v\mid u\in U_1, v\in U_2\}
	\end{align*}
	die Summe von $U_1$ und $U_2$. Gilt $U_1\cap U_2 = \{0\}$, so schreibt man $U_1\oplus U_2:= U_1+U_2$ und nennt es direkte Summe. Falls $U_1\oplus U_2=V$, dann nennt man $U_1$ und $U_2$ komplementär (nicht eindeutig).
\end{definition}
\begin{theorem}
	$V$ $K$-VR, $U_1,U_2\leq V$. Dann $U_1+U_2$ der kleinste UVR von $V$, der $U_1$ und $U_2$ enthält, d.h. $U_1+U_2=\mathrm{Lin}_V(U_1\cup U_2)$.
\end{theorem}
\prove beidseitige Teilmenge 
\begin{theorem}{Dimensionssatz}
	Sei $V$ $K$-VR. Sind $U_1, U_2 \leq V$ endlich dimensional, dann gilt\begin{align*}
		\dim_K(U_1+U_2)=\dim_K(U_1)+\dim_K(U_2)-\dim_K(U_1\cap U_2)
	\end{align*}
\end{theorem}
\prove Basis vom Schnitt erweitern, lineare Unabhängigkeit von extra Vektoren die nicht im Schnitt sind zeigen
















